\section{Introduction}

\subsection{Context and Motivation}
Techniques developed through game-playing AI have subsequently been applied and led to improvements in non-game domains including computer chip design and data centre optimisation\cite{generaladvancements}. Games provide controlled experimental environments because their rules and outcomes are clearly defined, allowing different AI approaches to be developed, tested and compared under consistent conditions\cite{yannakakis2025artificial}. They also provide opportunities to investigate human-AI interaction. Because games are widely played and designed to engage participants through continuous interaction, they offer an accessible context for studying how people respond to artificial players\cite{humanaiinteraction}.\\

Although games offer controlled environments and may have simple rules, they are not necessarily easy for AI to play. Each decision in a game may present several possible options, producing a rapidly expanding number of future states. Examining every possible move sequence can therefore become impractical. Games can also vary in length, scoring structure and sources of randomness, creating different demands for an AI agent\cite{yannakakis2025artificial}. These characteristics of a game can have an impact as to which approaches are suitable and how effectively they perform. Furthermore, available processing power and permitted computation time impose limits on an agent's playing strength\cite{deepblue}.\\

Traditional games such as Chess and Go have been studied extensively within game-playing AI, supported by long competitive histories, established strategic knowledge and substantial collections of recorded play\cite{alphago}. Research has consequently concentrated heavily on traditional board games, particularly deterministic, perfect-information games, while modern board games have received comparatively less attention\cite{tagframework}. These newer games can introduce additional challenges through randomness, hidden information, changing rules, varied player roles and different forms of interaction\cite{tagframework}. Applying established AI approaches to such environments may produce different results and reveal limitations that are less apparent in traditional games. It may also suggest new directions for game-playing AI research. \\

The project considers three broad approaches to game-playing AI with each approach selecting actions differently. Traditional search-based methods explore possible moves and their consequences repeatedly to identify promising moves. Simulation-based search methods estimate the value of moves by repeatedly simulating how the game might progress, while learning-based methods improve their strategies by gaining experience from repeated play\cite{yannakakis2025artificial}. Comparing these approaches across different game environments can help identify their respective strengths and limitations, as well as assess their suitability for the challenges presented by modern board games.\\


\subsection{Project Aim and Scope}
The project investigates how different game-playing AI approaches perform when applied to a small traditional game and a modern board game. Tic-tac-toe is chosen as the traditional game since it is widely understood, has simple rules and has been solved, which means that optimal play by both players will always result in a draw\cite{tictactoesolved}. The small state space allows an easier environment to introduce and validate the AI agents, while established strategies and existing implementations provide useful reference points for comparison\cite{tictactoeexisting}. Azul serves as the modern board game and holds additional challenges due to its longer game, wider range of possible actions, scoring system, round-based gameplay and a source of randomness due to its tile factory system\cite{azulrulebook}. The project focuses on the two-player version of Azul in order to maintain a consistent two-agent environment across both games while also avoiding the strategic complexity introduced by including more players.\\

Minimax is chosen as the algorithm for the search-based method, while Monte Carlo Tree Search is used as the simulation-based search method and Reinforcement Learning forms a broader learning-based category\cite{yannakakis2025artificial}. Within RL, the project investigates a progression from Tabular Q-Learning to Deep Q-Networks, and later an alternate RL approach using the policy-optimisation method Proximal Policy Optimisation\cite{yannakakis2025artificial}. Comparing these approaches are useful since they place their computation demands at different stages. Minimax and MCTS build up search trees to find the ideal action, although MCTS is able to return a result based on an allocated time budget\cite{mctshistory}. In contrast, the RL methods perform most of their computation during a separate training, allowing the agent to perform actions at a lower computational cost during play\cite{yannakakis2025artificial}. \\

The produced agents will be compared across several dimensions rather than relying only on playing strength. Playing strength will be determined through win rates against other agents, as well as game scores for Azul. Efficiency will also be considered using decision time, computation requirements, and training costs for RL agents. The comparison will distinguish between computation performed during training and play. Scalability and observed behaviour will also be investigated to identify how performance varies between games and what general strategies are found. These factors will assist in establishing the strengths, limitations and potential suitability against human opponents for each agent. \\ 